%\newpage
\section*{Conclusion}
This study successfully demonstrated the effectiveness of fuzzy logic and Iterative Learning Control (ILC) for three representative systems: inverted pendulum balancing, water tank level regulation, and repetitive trajectory tracking in discrete-time plants. A Mamdani fuzzy controller with 81 rules stabilized the inherently unstable inverted pendulum while accurately tracking the cart position with acceptable overshoot (16.7\%) and rapid settling. For the nonlinear water tank, the zero-order Sugeno FIS dynamically tuned $K_P$ and $K_I$ gains, enabling the fuzzy-corrected PID controller to achieve faster rise time, reduced overshoot, smoother transients, and negligible steady-state error across all setpoints ($SV_1=30$, $SV_2=60$, $SV_1+SV_2=90$) compared to conventional PI control. In the ILC section, a simple P-type scheme achieved excellent tracking for both linear and nonlinear discrete plants after only 60 iterations, with the Mean Squared Error converging rapidly to near zero. Overall, the results confirm that fuzzy logic provides a robust, model-free solution for nonlinear systems, while ILC effectively exploits task repetition to eliminate tracking errors. These intelligent techniques significantly enhance stability, tracking accuracy, and adaptability without requiring precise mathematical models, offering a solid foundation for future hardware implementation and advanced control applications.
